%% kapa: exam=1 round=32 year=2021 period=1 form=A subject=economics subject_ko=경제학원론 seq=7 number=47
%% source: Q-Net 감정평가사 기출문제 (cst00309) artlSeq=5208397 file=extracted/1008/5208397/A형/2021년도 제32회 감정평가사 1차 1교시 A형.pdf
%% conversion: kordoc 4.24.51 -> tools/gen1.py (auto); review=proofread-sonnet; flags=none
\begin{question}{47}
시장 수요이론에 관한 설명으로 옳지 않은 것을 모두 고른 것은?
\begin{framedtext}
ㄱ. 네트워크효과가 있는 경우 시장수요곡선은 개별 수요곡선의 수평 합이다. \par
ㄴ. 상품 소비자의 수가 증가함에 따라 그 상품 수요가 증가하는 효과를 속물효과(snob effect)라고 한다. \par
ㄷ. 열등재라도 대체효과의 절대적 크기가 소득효과의 절대적 크기보다 크면 수요곡선은 우하향한다. \par
ㄹ. 소득이 증가할 때 소비가 증가하는 재화는 정상재이다. \par
\end{framedtext}
\begin{choices}
  \item ㄱ, ㄴ
  \item ㄱ, ㄷ
  \item ㄱ, ㄹ
  \item ㄴ, ㄷ
  \item ㄴ, ㄹ
\end{choices}
\end{question}
\answer{1}
